% Label the appendix chapters "Appendix" in the table of contents, by redefining
% \chaptername inside the .toc from this point on. The line must live INSIDE the
% first included appendix: a write issued between two \include's migrates into the
% next include's .aux and would land in the .toc after this chapter's own entry.
\addtocontents{toc}{\string\def\string\chaptername{Appendix}}

\chapter{Reproducibility}
\label{app:first}

The appendices run at full width: the margin column is folded into the text block
after \verb|\appendix|, so tables and listings get 164.6\,mm instead of 115\,mm.
A wide table is therefore unnecessary here, and a normal table already spans the block.

\section{Configuration}
\label{app:first-config}

\begin{table}[htbp]
  \centering
  \small
  \caption[A full-width appendix table]{An appendix table. It uses the full 164.6\,mm, which is why the appendices carry no margin elements.}
  \label{tab:app-config}
  \begin{tabular}{@{}l p{6cm} p{5cm}@{}}
    \toprule
    Field & Value & Note \\
    \midrule
    Model & the model identifier & as published \\
    Revision & the commit of the model repository & pin it, tags move \\
    Precision & bfloat16, no quantization & \\
    Seeds & 42, 43, 44 & three per configuration \\
    \bottomrule
  \end{tabular}
\end{table}

\Cref{tab:app-config} prints as "Table A.1", and \Cref{app:first-config} prints as "Appendix A.1", because \verb|\crefalias| is set after \verb|\appendix|.
